% Layout adapted from the research memo shown in Lecture 1, sized for A4.
\documentclass[a4paper,10pt]{article}
\usepackage[margin=18mm,footskip=9mm]{geometry}
\usepackage{iftex}
\ifPDFTeX
  \usepackage[T1]{fontenc}
  \usepackage[utf8]{inputenc}
  \usepackage[scaled=0.96]{helvet}
  \renewcommand{\familydefault}{\sfdefault}
\else
  \usepackage{fontspec}
  \setmainfont{texgyreheros}[Extension=.otf,
    UprightFont=*-regular, BoldFont=*-bold,
    ItalicFont=*-italic, BoldItalicFont=*-bolditalic]
\fi
% For Chinese, set this switch true and compile with --engine xelatex.
\newif\ifChineseMemo
\ChineseMemofalse
\ifChineseMemo
  \ifPDFTeX
    \PackageError{finance-memo}{Chinese text requires XeLaTeX or LuaLaTeX}{}
  \else
    \usepackage[UTF8,fontset=fandol]{ctex}
    \renewcommand{\familydefault}{\sfdefault}
  \fi
\fi
\usepackage{booktabs,graphicx,xcolor,caption,xurl,fancyhdr}
\definecolor{MemoBlue}{HTML}{203E70}
\definecolor{MemoRule}{HTML}{B9C8DF}
\usepackage[colorlinks=true,linkcolor=MemoBlue,urlcolor=MemoBlue]{hyperref}
\captionsetup{font=small,labelfont=bf,skip=4pt}
\setlength{\parindent}{0pt}
\setlength{\parskip}{5pt}
\pagestyle{fancy}
\fancyhf{}
\renewcommand{\headrulewidth}{0pt}
\renewcommand{\footrulewidth}{0.3pt}

\newcommand{\ReportTitle}{[Company, portfolio, or decision]}
\newcommand{\ReportType}{RESEARCH MEMO}
\newcommand{\ReportDate}{[Report date]}
\newcommand{\DataDate}{[Data cutoff]}
\newcommand{\ReportScope}{[Period and comparison basis]}
\newcommand{\TeamName}{[Team]}
\newcommand{\DeskName}{[Business area]}
\newcommand{\ReportAuthor}{[Analyst]}
\newcommand{\Coverage}{[Coverage]}
\newcommand{\KeyFinding}{[One sentence stating the main supported finding]}
\fancyfoot[L]{\scriptsize\TeamName}
\fancyfoot[R]{\scriptsize\ReportType}

\begin{document}
\begin{minipage}[t]{0.70\textwidth}
\vspace{0pt}
{\small\ReportDate}\par\medskip
{\color{MemoBlue}\hrule height 0.5pt}\smallskip
{\large\ReportType}\par
{\LARGE\bfseries\ReportTitle}\par\medskip
\ReportScope
\end{minipage}\hfill
\begin{minipage}[t]{0.23\textwidth}
\vspace{0pt}\raggedleft
\colorbox{MemoBlue}{\parbox[c][19mm][c]{0.86\linewidth}{\centering\color{white}\bfseries\TeamName}}\par\medskip
\DeskName
\end{minipage}
\par\bigskip
{\color{MemoBlue}\hrule height 0.5pt}\medskip
{\large\KeyFinding}\par\bigskip

\begin{minipage}[t]{0.61\textwidth}
\vspace{0pt}\raggedright
\setlength{\parskip}{4pt}
\textbf{Scope}\par
[State the question, entities, reporting periods, and definitions.]

\textbf{Findings}\par
[Explain the verified results and their relevance to the reader's decision.]

\textbf{Interpretation}\par
[Separate reported facts, calculations, assumptions, and judgments.]

\textbf{Risks and limitations}\par
[State material data gaps, period mismatches, and uncertainty.]
\end{minipage}\hfill
\begin{minipage}[t]{0.36\textwidth}
\vspace{0pt}\small
% Author these non-floating wrappers separately. Use \captionof{table/figure},
% then a description and \input/\includegraphics of generated content.
\IfFileExists{sections/tables.tex}{\input{sections/tables.tex}}{
  \textbf{Key data}\par
  [Insert a compact checked table, including units and periods.]
}
\medskip
\IfFileExists{sections/figures.tex}{\input{sections/figures.tex}}{
  \textbf{Supporting chart}\par
  [Insert a chart with readable labels, units, and a source note.]
}
\end{minipage}
\par\vfill

{\color{MemoRule}\hrule height 0.4pt}\medskip
\begin{minipage}[t]{0.36\textwidth}
\vspace{0pt}\small
\setlength{\parskip}{3pt}
\textbf{PREPARED BY}\par\ReportAuthor\par\TeamName\par
\medskip\textbf{COVERAGE}\par\Coverage
\end{minipage}\hfill
\begin{minipage}[t]{0.60\textwidth}
\vspace{0pt}\footnotesize
\setlength{\parskip}{3pt}
\textbf{Sources and dates.} Data as of \DataDate.
[Identify the provider, retrieval date, original files, and source locators.]

\textbf{Checks.} [Record completed verification and any outstanding checks.]

\textbf{Preparation.} [Describe actual methods and relevant AI assistance.]
\end{minipage}
\end{document}
